\chapter{Methodology}
\label{ch:methodology}

Our simulation framework models a stylized market environment in which agents
repeatedly allocate fractional wealth.  This allows us to test theoretical
trade-offs under rigorous, repeatable conditions with full control over
distributional assumptions, horizon length, and regime structure.

% ---------------------------------------------------------------
\section{The Stylized Market Environment}
% ---------------------------------------------------------------

We model the market as the filtered probability space
$(\Omega,\Filt,\{\Filt_t\},\Prob)$ of Definition~\ref{def:filtered_prob_space},
with returns $r_t \mid S_t$ drawn from one of $M$ regime-specific distributions.

\subsection{Regime Configurations}

Three synthetic environments are used throughout the experiments:

\paragraph{Stationary.}
A single regime with $r_t \sim \Normal(\mu,\sigma^2)$, $\mu=0.08$,
$\sigma=0.15$.  Tests pure exploitation and validates the SLLN result of
Theorem~\ref{thm:kelly_slln}.

\paragraph{Adversarial Shock.}
A deterministic structural break at $T_{\text{shock}} = T/2$.  Bull market
$r_t \sim \Normal(0.08,\,0.15^2)$ for $t \leq T_{\text{shock}}$; Bear market
$r_t \sim t_3(-0.10,\,0.30^2)$ for $t > T_{\text{shock}}$.  The Student-$t$
with 3 degrees of freedom produces heavy-tailed returns that replicate the
Kelly-Ruin Paradox of Section~\ref{sec:kelly_ruin}.

\paragraph{Slow Gaussian Drift.}
The mean drifts linearly: $\mu_t = 0.08 + (-0.05 - 0.08)\cdot t/T$, with
$\sigma=0.15$ constant.  Tests agents' ability to unlearn a prior gradually
rather than responding to a single abrupt shock.

\subsection{Common Random Numbers (CRN)}
\label{sec:crn}

All agents in a given experiment are evaluated on the \emph{same} realized
return matrix $\mathbf{R} \in \R^{N_{\text{sim}}\times T}$, generated once
before any agent acts.  The CRN variance-reduction technique (L'Ecuyer, 1994)
ensures that differences in terminal wealth between agents are attributable
solely to strategy differences, not Monte Carlo sampling noise.

\begin{remark}
Any agent comparison that generates independent return paths for each agent
is statistically invalid under CRN.  All simulation loops in this codebase
share a single pre-generated \texttt{returns\_matrix} array.
\end{remark}

% ---------------------------------------------------------------
\section{Decision Agents}
\label{sec:agents}
% ---------------------------------------------------------------

\subsection{Kelly Oracle (Benchmark)}

The Oracle knows the true distribution parameters at every step and bets
$f^* = \mu/\sigma^2$ (Equation~\eqref{eq:gaussian_kelly}).  It represents the
theoretical upper bound and is used to compute the \emph{Oracle Regret}
$\mathcal{R}_t^{\text{oracle}} = \log W_t^{(f^*)} - \log W_t^{(\text{agent})}$
plotted in the regret dynamics figures.

\subsection{Naive Bayes Kelly}

Maintains Beta posteriors $(\alpha_k,\beta_k)$ and bets the posterior-mean
Kelly fraction $f^*(\hat p_k)$.  As formalized in
Proposition~\ref{prop:sticky_prior}, this agent suffers $O(T_0)$ detection lag
after a regime change.

\subsection{Thompson Sampling Kelly}

Samples $\tilde p_k \sim \mathrm{Beta}(\alpha_k,\beta_k)$ and bets
$f^*(\tilde p_k)$.  Bayesian posterior variance acts as implicit risk aversion:
high variance $\Rightarrow$ dispersed samples $\Rightarrow$ lower average bets
$\Rightarrow$ lower early-stage ruin risk versus the Naive plug-in.

\subsection{UCB Agent}

Applies the UCB1 rule of Theorem~\ref{thm:ucb_regret} to estimate each arm's
win probability, then maps the optimistic estimate through the Kelly formula.

\subsection{EXP3 Agent (Fractional Variant)}

Uses exponential-weight updates $w_k \leftarrow w_k\exp(\gamma\hat r_k/p_k)$
and allocates fractions proportional to the resulting mixing probabilities.
See Remark~\ref{rem:exp3_variant} for the deviation from canonical EXP3.

\subsection{Proposed Method 1: Volatility-Augmented HMM (Vol-HMM)}
\label{sec:vol_hmm}

The Vol-HMM addresses the Sticky Prior Paradox through three mechanisms.

\paragraph{Two-dimensional observation.}
At each step the agent observes $(r_t,v_t)$ where
$v_t = \hat\sigma_{[t-W,t]}$ is the rolling standard deviation over window
$W=5$.  Volatility spikes immediately upon a regime shift, providing an
early-warning signal the univariate return cannot.

\paragraph{Log-space forward algorithm.}
Regime beliefs are maintained via Equation~\eqref{eq:forward_logspace}.  The
joint emission log-likelihood is
\begin{equation}
    \ln b_j(r_t,v_t) \;=\;
    \ln\Normal(r_t;\,\mu_j,\,\sigma_j^2)
    \;+\;
    \ln\Gam(v_t;\,k_j,\,\theta_j),
    \label{eq:vol_hmm_emission}
\end{equation}
where emission parameters are specified a priori
(Remark~\ref{rem:hmm_specified}):
$(\mu_{\text{bull}},\sigma_{\text{bull}},k_{\text{bull}},\theta_{\text{bull}})
 = (0.08,\,0.15,\,2.0,\,0.075)$ and
$(\mu_{\text{bear}},\sigma_{\text{bear}},k_{\text{bear}},\theta_{\text{bear}})
 = (-0.10,\,0.30,\,4.0,\,0.15)$.

\paragraph{CUSUM change-point trigger.}
A CUSUM statistic accumulates evidence of a negative mean shift:
\begin{equation}
    S_t^- \;=\; \max\!\bigl(0,\;
    S_{t-1}^- - z_t - k_{\mathrm{drift}}\bigr),
    \qquad
    z_t = \frac{r_t - \mu_{\text{bull}}}{\sigma_{\text{bull}}},
    \label{eq:cusum}
\end{equation}
with $k_{\mathrm{drift}}=0.5$.  When $S_t^- > h=3.0$, a $+2$ log-likelihood
bonus is added to the bear-state emission, accelerating belief reallocation.

\paragraph{Transition constraint.}
$A_{ii} \leq 0.95$ prevents sticky-prior re-emergence through the transition
dynamics.

\paragraph{Bet sizing.}
$f_t = \Prob(\text{bull}\mid\Filt_t)\cdot f^*_{\text{bull}}$, reducing
exposure linearly as belief mass flows to the bear state.

\subsection{Proposed Method 2: Risk-Constrained Kelly (CPPI)}
\label{sec:cppi}

\begin{definition}[CPPI Floor and Cushion]
\label{def:cppi}
Let $\bar W_t = \max_{s\leq t}W_s$ be the running peak wealth.  The
\emph{floor} is $F_t = (1-D_{\max})\bar W_t$ and the
\emph{cushion} is $C_t = W_t - F_t$.
\end{definition}

\begin{proposition}[CPPI Drawdown Bound]
\label{prop:cppi_bound}
Under the leverage rule $\lambda_t = \min(1,\, m\cdot C_t / W_t)$ with
multiplier $m > 0$, the wealth satisfies
\begin{equation}
    W_{t+1} \;\geq\; F_t\bigl(1 - m\,C_t\,|r_{\min}|/W_t\bigr),
    \label{eq:cppi_bound}
\end{equation}
where $r_{\min}$ is the worst single-period return.  In particular, when
$m\,|r_{\min}| \leq 1$, the drawdown from peak never exceeds $D_{\max}$.
\end{proposition}

\begin{proof}
The wealth update is
$W_{t+1} = W_t(1 + \lambda_t r_t) = W_t + m C_t r_t$ (ignoring the min cap).
In the worst case $r_t = r_{\min} < 0$:
$W_{t+1} = W_t + m C_t r_{\min} = F_t + C_t + m C_t r_{\min}
= F_t + C_t(1 + m r_{\min})$.
This is $\geq F_t$ iff $1 + m r_{\min} \geq 0$, i.e., $m|r_{\min}| \leq 1$.
In continuous time this condition is always satisfiable; in discrete time a
single return $r_t < -1/m$ can breach the floor (documented
experimentally in Chapter~\ref{ch:experiments}).
\end{proof}

\begin{remark}[Cost of Survival]
Proposition~\ref{prop:cppi_bound} makes precise the trade-off between safety
and growth.  A tighter floor ($D_{\max}$ small) requires a smaller cushion to
leverage, which forces $\lambda_t$ toward zero whenever $W_t$ is close to its
peak.  This is the ``cost of survival'': lower average leverage implies lower
expected log-growth than unconstrained Kelly, but the floor provides the
insurance against catastrophic loss.
\end{remark}

% ---------------------------------------------------------------
\section{Market Friction and Transaction Costs}
% ---------------------------------------------------------------

To bridge theoretical growth and empirical deployment, we model transaction
costs at each step.  Given betting fraction $f_t$ and drift-adjusted previous
fraction $\hat f_{t-1} = f_{t-1}(1+r_{t-1})/(1+f_{t-1}r_{t-1})$:
\begin{equation}
    W_{t+1} \;=\; W_t\cdot\bigl(1 + f_t r_t - c\,|f_t - \hat f_{t-1}|\bigr),
    \label{eq:friction}
\end{equation}
where $c$ is the friction in basis points ($\times 10^{-4}$).  Tested values:
$c\in\{0,10,25,50,100\}$ bps.  High-turnover agents (EXP3, UCB) rebalance
aggressively and therefore incur proportionally higher fees.

% ---------------------------------------------------------------
\section{Evaluation Metrics}
% ---------------------------------------------------------------

Strategies are evaluated on four complementary objectives:

\begin{itemize}
    \item \textbf{Median Terminal Wealth} $\tilde W_T$: the 50th percentile
    across $N_{\text{sim}}$ paths.  The median is used because the terminal
    wealth distribution is log-normal with a heavy right tail; the mean is
    dominated by rare extreme-growth paths.

    \item \textbf{Probability of Ruin}
    $\hat\rho = N_{\text{ruined}}/N_{\text{sim}}$.
    Statistical significance between two agents is assessed by Fisher's
    Exact Test on the $2\times 2$ ruin/survival contingency table.

    \item \textbf{Annualized Sharpe Ratio}
    $\mathcal{S} = (\bar R/\hat\sigma_R)\sqrt{252}$,
    computed via the $O(T)$ running-statistics recursion
    $\bar R_t = \bar R_{t-1} + (R_t - \bar R_{t-1})/t$.

    \item \textbf{Annualized Sortino Ratio}
    $\mathcal{SOR} = (\bar R/\hat\sigma_\downarrow)\sqrt{252}$,
    where $\hat\sigma_\downarrow^2 = T^{-1}\sum_t\min(R_t,0)^2$.
\end{itemize}

The \emph{Oracle Regret} $\log W_t^{(f^*)} - \log W_t^{(\text{agent})}$ is
plotted over time to reveal the dynamic convergence rate of each agent
relative to the growth-optimal benchmark.
